\section{Objetivos}

\subsection{Objetivo general}

Diseñar, implementar y evaluar una plataforma local de orquestación de agentes
de lenguaje que adapte la granularidad de un objetivo de software a su forma
semántica y entregue únicamente resultados verificables sobre un commit exacto.

\subsection{Objetivos específicos}

\begin{enumerate}
  \item Separar el juicio semántico que propone un corte de la política
  determinista que decide si aceptarlo y del compilador que lo transforma en un
  grafo ejecutable.
  \item Representar jerarquía, disponibilidad de artefactos, compatibilidad de
  interfaces y riesgo de conflicto mediante relaciones distintas y tipadas.
  \item Aislar cada intento en un worktree, controlar su alcance y adoptar el
  resultado sólo contra sus entradas exactas.
  \item Vincular cada criterio de aceptación con evidencia obtenida sobre el
  candidate SHA, incluyendo baseline, controles negativos e integridad de
  tests.
  \item Integrar de abajo hacia arriba y publicar sólo después de validar el
  candidato final y emitir un recibo de entrega.
  \item Persistir la historia del run como eventos append-only, separando hechos
  de dominio, snapshots y trazas diagnósticas.
  \item Evaluar dos proposiciones acotadas mediante una serie pre-registrada,
  repetida y verificada por un oráculo externo.
\end{enumerate}

\section{Preguntas de investigación e hipótesis}

\begin{description}[leftmargin=2.2cm,style=nextline]
  \item[\textbf{PI-1}] ¿Puede ManyHands llevar cambios de distinta forma desde
  el objetivo hasta una entrega confirmada cuyo commit exacto pase un oráculo
  externo?
  \item[\textbf{PI-2}] ¿Selecciona la política una topología multi-hoja para una
  tarea que atraviesa fronteras de integración y una única hoja para una tarea
  cohesiva?
\end{description}

PI-1 se operacionaliza como H-F1: cuatro celdas \texttt{completed} y
\texttt{delivered}, cuatro SHAs no vacíos y cuatro oráculos PASS. PI-2 se
operacionaliza como H-F2: dos topologías de tres hojas para la tarea multi-capa
y dos topologías de una hoja para la tarea cohesiva. Las reglas completas se
presentan en el capítulo~\ref{cap:evaluacion}.

Estas preguntas no incluyen superioridad frente a no descomponer. Tampoco
preguntan si la topología observada es globalmente óptima. El diseño final elige
proposiciones que el instrumento puede medir de forma completa, en vez de
mantener hipótesis históricas que quedaron fuera del régimen o no produjeron
candidatos.

\section{Alcance y limitaciones}

\sistema\ es un plano de control local y de un solo usuario. Planning,
scheduling, persistencia, integración y validación se ejecutan en la máquina del
desarrollador; la inferencia se obtiene mediante CLI de proveedores remotos. No
se afirma privacidad absoluta del código enviado al modelo.

La evaluación confirmatoria usa un target Node ESM pequeño, un modelo, una
política y dos repeticiones por forma de tarea. Permite estudiar viabilidad y
trazabilidad end-to-end, pero no escalabilidad, superioridad, causalidad,
optimalidad ni generalización estadística.

Las series históricas no se borran: Warehouse queda 1/8, G5 negativo, G6
inconcluso, G7 sin candidatos y SP2 como piloto. Se reportan como antecedentes
adversos o formativos, pero no se combinan con las cuatro celdas finales.

Quedan fuera del alcance: selección adaptativa de modelos, soporte multiusuario,
aislamiento por contenedores, calibración empírica de pesos, generalización a
otros lenguajes y auditoría externa de accesibilidad. También se declara como
gap de transición la materialización selectiva de artefactos por archivo: el
contrato puede expresarla, pero el driver actual aún transporta el commit
producido completo.
